\setcounter{table}{0}
\renewcommand{\thetable}{A-\arabic{table}}

\section{附录 A　高分二号（GF-2）卫星技术参数}

本附录汇总高分二号（GF-2）卫星的公开技术参数，作为正文第七章“项目技术的支撑”的资料性补充。正文侧重说明各参数与本项目监测需求的对应关系，本附录则给出完整的技术清单，便于查阅与核对。表中参数依据该卫星官方公布数据整理。

\subsection{A.1 平台与载荷主要参数}

\begin{table}[htbp]
    \centering
    \caption{高分二号（GF-2）卫星平台与载荷主要参数}
    \label{tab:gf2-platform}
    \small
    \begin{tabular}{|p{3.4cm}|p{5.4cm}|p{4.6cm}|}
        \hline
        \textbf{参数项} & \textbf{技术指标} & \textbf{说明} \\
        \hline
        发射时间 & 2014 年 8 月 19 日 & 于太原卫星发射中心发射 \\
        \hline
        轨道类型 & 太阳同步回归轨道 & 降交点地方时约上午 10:30 \\
        \hline
        轨道高度 & 631 km & 兼顾覆盖范围与成像分辨率 \\
        \hline
        回归周期 & 69 天 & 不侧摆成像时的全球覆盖周期 \\
        \hline
        侧摆能力 & $\pm35^{\circ}$ & 具备快速姿态机动能力 \\
        \hline
        重访周期 & 侧摆条件下约 5 天 & 满足“周级”动态监测的频次需求 \\
        \hline
        全色分辨率 & 星下点 0.8 m & 我国首颗分辨率优于 1 m 的民用光学遥感卫星 \\
        \hline
        多光谱分辨率 & 星下点 3.2 m & 与全色影像融合后提升地类判别能力 \\
        \hline
        成像幅宽 & 45 km & 由两台相机拼接实现 \\
        \hline
        定位精度 & 无地面控制点时约 50 m（CE90） & 保障多时相影像的几何配准精度 \\
        \hline
    \end{tabular}
    \par\vspace{2pt}
    {\footnotesize 注：表中“星下点”指卫星垂直投影点处的成像分辨率，为载荷实际成像指标；相机标称指标为全色 1 m、多光谱 4 m。}
\end{table}

\subsection{A.2 谱段设置与主要用途}

\begin{table}[htbp]
    \centering
    \caption{高分二号（GF-2）卫星谱段设置及本项目用途}
    \label{tab:gf2-bands}
    \small
    \begin{tabular}{|p{2.4cm}|p{3.4cm}|p{7.6cm}|}
        \hline
        \textbf{谱段} & \textbf{波长范围} & \textbf{主要用途} \\
        \hline
        全色 & 0.45--0.90 $\mu$m & 提供高分辨率灰度影像，用于地块边界、田埂与侵蚀沟等几何细节判读 \\
        \hline
        蓝 & 0.45--0.52 $\mu$m & 水体识别与大气校正，辅助地表反射率反演 \\
        \hline
        绿 & 0.52--0.59 $\mu$m & 表征植被反射特征，用于作物长势与植被覆盖监测 \\
        \hline
        红 & 0.63--0.69 $\mu$m & 叶绿素吸收波段，与近红外组合可计算 NDVI 等植被指数 \\
        \hline
        近红外 & 0.77--0.89 $\mu$m & 对植被与土壤水分敏感，用于植被覆盖度与土壤有机质反演 \\
        \hline
    \end{tabular}
\end{table}

\subsection{A.3 本项目的数据使用方式}

在本项目中，上述影像按以下方式进入业务流程：首先对全色与多光谱影像进行辐射定标与大气校正，并以全色影像为基准完成几何配准，使两个时相的影像在同一空间基准上逐像元对应；随后进行全色与多光谱融合，在保留多光谱信息的同时将空间分辨率提升至全色水平；最后将融合影像统一裁剪为 $256 \times 256$ 像素的影像块，作为变化检测模型的输入。

四类多光谱谱段在其中承担不同职责：可见光波段（蓝、绿、红）用于地类判别与植被指数计算，近红外波段则提供植被与土壤水分信息，二者结合使模型不仅能判断“地表是否发生变化”，还能为变化类型的判别提供光谱依据。这一处理流程与正文所述的“数据采集—智能检测—分析决策”技术闭环相衔接，是其中数据环节的具体实现方式。
